\documentclass[11pt,a4paper]{article}

\usepackage[margin=1in]{geometry}
\usepackage{amsmath}
\usepackage{amssymb}
\usepackage{amsthm}
\usepackage{booktabs}
\usepackage{xcolor}
\usepackage{hyperref}

\hypersetup{colorlinks=true, linkcolor=blue!50!black, urlcolor=blue!50!black,
            citecolor=blue!50!black}

\newtheorem{proposition}{Proposition}
\newtheorem{definition}{Definition}

\DeclareMathOperator*{\argmax}{arg\,max}
\newcommand{\Ind}{\mathbb{I}}
\newcommand{\Prob}{\mathbb{P}}
\newcommand{\Exp}{\mathbb{E}}

\title{\textbf{Rank Inference from a Non-Transitive Pairwise Oracle}\\
\large Query Complexity, Parallel Span, and an Empirical Accuracy Ceiling}
\author{Working draft}
\date{30 July 2026}

\begin{document}
\maketitle

\begin{abstract}
We study the problem of ordering a finite set of texts by fitness to a
natural-language criterion, where the only available signal is a pairwise oracle
realized by a large language model. Such an oracle is noisy and, critically, need
not be transitive: the induced tournament may contain directed cycles, so no
total order is consistent with all observed comparisons and comparison-sorting is
ill-posed. We adopt a two-stage formulation that separates \emph{query
scheduling} from \emph{rank inference}: an adaptive schedule selects which pairs
to query under a budget, and a regularized Bradley--Terry estimator fits latent
scores to all observed outcomes, absorbing cycles by construction. Within this
formulation we establish four results. First, ranking error decreases by
approximately $65\%$ as the query budget grows from $\Theta(n\log n)$ to
$\binom{n}{2}$, with sharply diminishing returns, and adaptively targeted queries
improve on uniformly random ones by only ${\sim}8\%$ at equal budget --- density
dominates targeting. Second, the objective determines the algorithm: when the
goal is to \emph{select} a good $k$-subset rather than to order it, a screening
schedule attains half the selection regret of a full sort at $56\%$ of its query
cost. Third, we give a span analysis showing that a merge-sort schedule has
$\Theta(n)$ sequentially dependent queries whereas round-based schedules have
span $O(r)$ in the number of rounds, so the selection-oriented schedule is
simultaneously cheaper and asymptotically more parallel. Fourth, on $500$
sentences of the Stanford Sentiment Treebank we obtain Kendall $\tau_b = 0.613$
using $3.4\%$ of all pairs, and find that oracle accuracy increases
monotonically with ground-truth separation, from $53.7\%$ on near-ties to
$98.4\%$ on well-separated pairs. This establishes an accuracy ceiling
attributable to the oracle rather than to the inference procedure. The ceiling is
severe under exact rank recovery yet largely benign under a subset-quality
criterion: at $k=10$ a precision of $0.2$ coincides with a selection regret of
$0.235$, the returned subset capturing ${\sim}77\%$ of the available quality
margin.
\end{abstract}

\section{Introduction}

Let $X = \{x_1,\dots,x_n\}$ be a set of texts and let $p$ be a criterion
expressed in natural language. We wish to order $X$ by fitness to $p$, given
access only to an oracle that compares two elements at a time.

Realizing this oracle with a language model introduces two departures from the
classical comparison model. The oracle is \emph{stochastic}: repeated queries, or
queries with the two candidates exchanged, may disagree. More consequentially it
is \emph{intransitive}: the relation it induces may contain cycles
$x_a \succ x_b \succ x_c \succ x_a$. Comparison sorting presumes a total order;
applied to an intransitive relation its output is a function of the comparison
schedule and of implementation details rather than of the data. A third
constraint is economic: each query is a model inference, so the query count is
the dominant cost and $\binom{n}{2}$ queries are infeasible for even moderate
$n$.

We therefore treat comparisons as \emph{evidence} to be aggregated rather than
as an ordering to be realized. Section~\ref{sec:formulation} formalizes the query
model and the two objectives we consider; Section~\ref{sec:method} specifies the
estimator and the schedules and analyzes their query complexity and parallel
span; Section~\ref{sec:experiments} reports simulation and empirical results.

\section{Problem formulation}
\label{sec:formulation}

\subsection{Query model}

\begin{definition}[Comparison oracle]
An oracle $\mathcal{O}$ accepts an unordered pair $\{i,j\}$ and returns a verdict
in $\{i, j, \varnothing\}$, where $\varnothing$ denotes a declared tie. A
\emph{schedule} is an adaptive procedure that issues a sequence of queries, each
possibly depending on previous verdicts, and terminates after at most $B$
distinct queries.
\end{definition}

We impose the invariant that each unordered pair is queried at most once, its
verdict being reused thereafter. This is a modeling choice with a substantive
consequence: it renders the oracle \emph{deterministic} on its domain, so
stochasticity manifests only as a fixed but possibly erroneous verdict per pair.
An erroneous verdict cannot be corrected by resampling; it can only be
outweighed by other edges incident to the same vertices. We revisit this in
Section~\ref{sec:density}.

Verdicts induce a weighted digraph on $X$ with $w_{ij}$ counting verdicts in
favour of $i$ over $j$; a tie contributes $\tfrac12$ to each direction. This
digraph need not be acyclic.

To neutralize the oracle's known sensitivity to presentation
order~\cite{zheng2023}, each query is issued with the candidate order drawn
uniformly at random and the verdict re-oriented afterwards, so that no positional
preference is systematically attributed to either operand.

\subsection{Objectives}

Two objectives are of interest, and they are not interchangeable.

\paragraph{Full ranking.} Recover a permutation agreeing with a latent order. We
measure agreement by Kendall's $\tau_b$, which corrects for ties in the reference
ordering~\cite{kendall1945}:
\begin{equation}
\tau_b = \frac{C - D}{\sqrt{(C+D+T_x)(C+D+T_y)}},
\end{equation}
with $C$, $D$ the numbers of concordant and discordant pairs and $T_x$, $T_y$
the numbers of pairs tied in only one of the two orderings. In simulation we
also report the normalized Kendall distance, the fraction of pairs misordered.

\paragraph{Top-$k$ selection.} Return a subset $S \subseteq X$, $|S| = k$, of
high-quality elements, with no requirement that $S$ be internally ordered. Let
$q: X \to \mathbb{R}$ denote latent quality and $\bar q(A)$ the mean of $q$ over
$A$. We report both
\begin{equation}
\text{precision}(S) = \frac{|S \cap X_k^{*}|}{k},
\qquad
\mathrm{regret}(S) = \frac{\bar q(X_k^{*}) - \bar q(S)}
                          {\bar q(X_k^{*}) - \bar q(X)},
\label{eq:regret}
\end{equation}
where $X_k^{*}$ is an optimal $k$-subset. Regret vanishes for an optimal
selection and equals $1$ in expectation for a uniformly random one. The two
measures differ precisely on near-ties: exchanging an element of $X_k^{*}$ for
one of nearly equal quality reduces precision by $1/k$ while leaving regret
almost unchanged. Since near-ties are where the oracle is least reliable
(Section~\ref{sec:sst}), the choice of measure determines whether oracle error
is consequential.

\section{Method}
\label{sec:method}

\subsection{Rank inference}

We aggregate verdicts with the Bradley--Terry model~\cite{bradley1952}, under
which each item carries a strength $\gamma_i > 0$ and
$\Prob(i \succ j) = \gamma_i / (\gamma_i + \gamma_j)$. The model is fitted by the
Zermelo/MM iteration~\cite{zermelo1929,hunter2004}, which we regularize:
\begin{equation}
\gamma_i^{(t+1)} \;=\;
\frac{\alpha + \sum_{j \neq i} w_{ij}}
     {\dfrac{2\alpha}{\gamma_i^{(t)} + 1} \;+\;
      \sum_{j \neq i} \dfrac{w_{ij} + w_{ji}}{\gamma_i^{(t)} + \gamma_j^{(t)}}},
\label{eq:mm}
\end{equation}
followed by renormalization of $\{\gamma_i\}$ to unit geometric mean. The terms
in $\alpha$ constitute a Jeffreys-type prior contributing $\alpha$ fictitious
wins and $\alpha$ fictitious losses against a virtual opponent of unit strength;
we take $\alpha = \tfrac12$.

Regularization is not cosmetic. When the comparison digraph is linearly
separable --- which occurs whenever the observed verdicts admit a consistent
total order --- the unregularized likelihood has no interior maximum and the
extreme strengths diverge. On an $8$-element instance we observed fitted
strengths spanning $2.5\times10^{-3}$ to $1.46\times10^{4}$, rendering
magnitudes uninterpretable while leaving the induced order unchanged. The prior
yields a proper posterior mode and bounded strengths. We report the additive
parameterization
\begin{equation}
s_i = \log \gamma_i, \qquad \Prob(i \succ j) = \sigma(s_i - s_j),
\end{equation}
so that scores are centred near zero and differences read as log-odds.

Two properties of this estimator are relevant to intransitivity. Cycles are
absorbed rather than broken: the fit maximizes the likelihood of the entire
verdict set rather than requiring each edge to be satisfied, and on a symmetric
$3$-cycle it returns $s_a = s_b = s_c = 0$, reporting a tie. Second, strengths
are opponent-weighted, correcting a deficiency of raw win counts: on a
constructed $7$-element instance, an item with two victories over strong
opponents received $s = 2.11$ while an item with two victories over weak
opponents received $s = 0.73$, and the former was ranked above an item with
three victories.

Elements with identical text are merged before inference, since such a pair is
tied by definition and yields no information; they receive equal scores and
adjacent positions.

\subsection{Schedules}

We consider three families of schedule.

\begin{itemize}
  \item \textbf{Merge sort.} A merge sort issues $\Theta(n \log n)$ queries. Its
    output permutation is discarded; only the verdicts are retained. It serves
    purely as a query schedule.
  \item \textbf{Round-based (Swiss).} In each of $r$ rounds, items are
    partitioned into disjoint pairs and all pairs are queried. The first round
    pairs at random; subsequent rounds pair items adjacent in the current fitted
    order, concentrating queries where the ordering is contested. Cost
    ${\approx} rn/2$.
  \item \textbf{Screening for selection.} $r$ Swiss rounds, followed by
    exhaustive querying among the top $ck$ candidates of the fitted order, which
    settles the boundary of the retained set. Cost
    ${\approx} rn/2 + \binom{ck}{2}$.
  \end{itemize}

A refinement phase may additionally query pairs adjacent in the fitted order
that remain unqueried, refitting between batches.

\subsection{Query complexity and parallel span}
\label{sec:span}

Because queries are the dominant cost and may be issued concurrently, two
quantities matter: the total number of queries, and the \emph{span}, i.e.\ the
length of the longest chain of queries in which each depends on the outcome of
its predecessor. Span lower-bounds latency irrespective of available parallelism.

\begin{proposition}[Span of a merge-sort schedule]
\label{prop:span}
Merging two sorted sequences of lengths $\lceil n/2\rceil$ and
$\lfloor n/2\rfloor$ requires, in the worst case, $n-1$ queries each of which
depends on the outcome of its predecessor. Consequently the span $S(n)$ of the
merge-sort schedule satisfies $S(n) = S(n/2) + (n-1)$, giving
$S(n) = \Theta(n)$.
\end{proposition}

\begin{proof}[Proof sketch]
Within a merge, the choice of the pair compared at step $t+1$ is determined by
which operand was emitted at step $t$; hence the queries of a single merge form a
dependency chain of length up to $n-1$. Sibling subproblems are independent, so
the span recursion adds only the cost of the root merge to that of one
subproblem, and unrolling $S(n) = S(n/2) + (n-1)$ yields $\Theta(n)$.
\end{proof}

By contrast, the $\lfloor n/2 \rfloor$ queries of a Swiss round act on disjoint
pairs and are mutually independent, so a round contributes $1$ to the span and
the $r$-round schedule has span $r$; the screening schedule adds a single
parallel batch, giving span $r+1$. Thus for fixed $r$ the round-based schedules
have span $O(1)$ in $n$, against $\Theta(n)$ for merge sort --- an asymptotic
separation in attainable latency, independent of the constant-factor cost of an
individual query.

\section{Experiments}
\label{sec:experiments}

\subsection{Simulation methodology}

To evaluate against known ground truth we instantiate a synthetic oracle obeying
the model of Section~\ref{sec:method}. Item $i$ is assigned latent strength
$\theta_i$ on an evenly spaced grid spanning $\Delta$ logits, and the oracle
returns $i$ over $j$ with probability $\sigma(\theta_i - \theta_j)$, sampled once
per pair and thereafter fixed, in accordance with the query model. With
$n$ items spanning $\Delta$ logits, adjacent items differ by
$\Delta/(n-1)$ logits and the oracle is correspondingly close to
uninformative on them; we choose $\Delta = 4$, giving adjacent-pair accuracy of
$51.7\%$ at $n = 60$ and $50.5\%$ at $n = 200$. This deliberately reproduces the
regime observed empirically in Section~\ref{sec:sst}. All simulation figures
average over independent trials with distinct seeds.

\subsection{Query density versus accuracy}
\label{sec:density}

We fix $n = 60$ and $24$ trials and vary the query budget. The window-$w$
schedule augments a merge-sort base with queries between items within $w$
positions of one another in the current fitted order, refitting between rounds.
The final column applies the same number of additional queries to uniformly
random pairs, isolating the value of targeting from that of density.

\begin{table}[h]
\centering
\begin{tabular}{lrrrr}
\toprule
Schedule & Queries & Kendall dist. & Top-10 & Random @ equal budget \\
\midrule
Merge sort only      &  227 & 0.2293 & 5.5/10 & --- \\
Window $w=1$         &  328 & 0.1680 & 6.0/10 & 0.1813 \\
Window $w=2$         &  412 & 0.1464 & 6.8/10 & 0.1597 \\
Window $w=4$         &  558 & 0.1224 & 7.0/10 & 0.1361 \\
Window $w=8$         &  784 & 0.1044 & 7.4/10 & 0.1167 \\
Exhaustive           & ${\sim}1997$ & 0.0792 & 8.3/10 & 0.0743 \\
\bottomrule
\end{tabular}
\caption{Normalized Kendall distance (lower is better) versus query budget;
$n=60$, $24$ trials.}
\label{tab:density}
\end{table}

Error decreases monotonically in the budget, by ${\sim}65\%$ from
$\Theta(n\log n)$ to exhaustive, but with pronounced concavity: the first
${\sim}100$ additional queries ($+44\%$ cost) remove $27\%$ of the error, whereas
attaining the exhaustive figure costs a further $6\times$ for an additional
$53\%$. Targeting yields a consistent but modest advantage over random selection
at matched budget, ${\sim}8\%$ in relative error across all budget levels; the
dominant factor is density rather than the choice of which pairs to query.

This is consistent with the deterministic-oracle property noted in
Section~\ref{sec:formulation}: additional edges cannot resample an erroneous
verdict, and act instead by increasing the number of constraints incident to each
vertex, so that a single incorrect edge is outweighed. A residual Kendall
distance of $0.079$ persists under exhaustive querying, an error floor
determined by oracle quality alone.

\subsection{Objective choice: ranking versus selection}
\label{sec:topk}

We now fix the selection objective of Eq.~\eqref{eq:regret} with $n = 200$,
$k = 10$ and $20$ trials, and compare schedules.

\begin{table}[h]
\centering
\begin{tabular}{lrrr}
\toprule
Schedule & Queries & Precision & Normalized regret \\
\midrule
Merge sort (full order)     & 1042 & $24.0\%$ & 0.3740 \\
Swiss $r=2$                 &  200 & $6.5\%$  & 0.7298 \\
Swiss $r=4$                 &  400 & $15.5\%$ & 0.3622 \\
Swiss $r=6$                 &  598 & $23.5\%$ & 0.2556 \\
Screening $r=4$, $c=2$      &  582 & $26.0\%$ & \textbf{0.1866} \\
Screening $r=4$, $c=4$      & 1152 & $42.0\%$ & \textbf{0.1012} \\
\bottomrule
\end{tabular}
\caption{Top-$k$ selection quality; $n=200$, $k=10$, $20$ trials.}
\label{tab:topk}
\end{table}

Two conclusions follow. First, the two measures of Eq.~\eqref{eq:regret} rank the
same output very differently. The merge-sort schedule attains precision $0.24$,
yet its regret of $0.374$ indicates that the returned subset captures
${\sim}63\%$ of the quality margin available over a random selection. Precision
is the more pessimistic measure by construction whenever the upper region of the
list is tightly clustered, since it charges full cost for exchanging elements of
nearly equal quality.

Second, a full sort is dominated for this objective. Screening with $c = 2$
attains half the regret of the merge-sort schedule using $56\%$ of the queries,
and exceeds it on precision as well; Swiss with $r = 4$ matches its regret at
$38\%$ of the query cost. The full sort expends budget on two tasks the objective
does not reward: ordering elements that will be excluded from $S$, and ordering
elements within $S$. Combined with Proposition~\ref{prop:span}, the screening
schedule is thus preferable in queries and in span simultaneously.

\subsection{Confidence signals}
\label{sec:conf}

The $8\%$ margin available to targeting (Section~\ref{sec:density}) motivates a
per-query uncertainty estimate with which to prioritize pairs. We compare two
estimators obtainable from the same oracle: a verbalized confidence, elicited by
requesting a probability alongside the verdict, and the normalized
log-probability assigned to the verdict token. Three pairs of differing
difficulty were each queried in both presentation orders; a well-behaved
estimator should invert its verdict label when the operands are exchanged.

\begin{table}[h]
\centering
\begin{tabular}{lll}
\toprule
Pair difficulty & Verbalized & Token log-probability \\
\midrule
Well separated & $1.00 / 1.00$, order-stable  & $0.999 / 0.998$, order-stable \\
Ambiguous      & $0.60 / 0.50$, order-flipped & $0.608 / 0.642$, order-stable \\
Near-tie       & $0.60 / 0.60$, order-stable  & $0.525 / 0.825$, order-flipped \\
\bottomrule
\end{tabular}
\caption{Uncertainty estimators on pairs of graded difficulty.}
\label{tab:conf}
\end{table}

Both estimators separate well-separated pairs (${\approx}1.0$) from difficult
ones ($0.5$--$0.65$), which suffices for prioritization. The verbalized estimator
is however severely quantized, taking only three distinct values over six
queries and concentrating on round numbers, consistent with reported
miscalibration of elicited confidence~\cite{tian2023}; the log-probability
estimator is continuous and requires no additional generated tokens. Neither is
order-invariant, each inverting on one difficult pair, which supports the
symmetrization of Section~\ref{sec:formulation} and suggests that agreement
across both presentation orders is the more reliable estimator, at twice the
query cost.

\subsection{Empirical evaluation on the Stanford Sentiment Treebank}
\label{sec:sst}

We evaluate against SST~\cite{socher2013}, selected because its labels are
continuous on $[0,1]$ rather than categorical; datasets with few graded
relevance levels induce large numbers of reference ties and render $\tau_b$
uninformative. Joining sentences to phrase-level labels yields $11{,}286$
sentences with real-valued scores spanning the full interval over $82$ distinct
levels, with median length $16$ words; $569$ sentences were discarded for lack of
an exact label match. Sentences are converted from Penn Treebank tokenization to
orthographic form. The oracle is a locally served $9$B-parameter instruction-tuned
model under greedy decoding, and the criterion is expressed as ``the most
positive and enthusiastic sentiment about the film''.

\begin{table}[h]
\centering
\begin{tabular}{lrrrr}
\toprule
$n$ & Queries & Fraction of $\binom{n}{2}$ & $\tau_b$ & Spearman $\rho$ \\
\midrule
$40$ (distinct labels) &  200 & $25.6\%$ & 0.700 & 0.890 \\
$500$                  & 4227 & $3.4\%$  & 0.613 & 0.810 \\
\bottomrule
\end{tabular}
\caption{Rank recovery on SST.}
\label{tab:sst}
\end{table}

\begin{table}[h]
\centering
\begin{tabular}{lrr}
\toprule
Ground-truth separation & Accuracy ($n=40$) & Accuracy ($n=500$) \\
\midrule
$0.00$--$0.05$ & $59.2\%$  & $53.7\%$ \; ($6474/12048$) \\
$0.05$--$0.15$ & $60.6\%$  & $64.5\%$ \; ($16471/25536$) \\
$0.15$--$0.30$ & $86.2\%$  & $79.5\%$ \; ($23708/29831$) \\
$0.30$--$0.60$ & $100.0\%$ & $94.3\%$ \; ($40933/43424$) \\
$0.60$--$1.00$ & $100.0\%$ & $98.4\%$ \; ($11674/11869$) \\
\bottomrule
\end{tabular}
\caption{Oracle accuracy as a function of the ground-truth label gap.}
\label{tab:gap}
\end{table}

Table~\ref{tab:gap} is the principal empirical result. Oracle accuracy increases
monotonically with ground-truth separation, from near chance on near-ties to
${\sim}98\%$ on well-separated pairs, and the profile replicates across a
$60$-fold change in the number of judged pairs. The accuracy ceiling is therefore
a property of the oracle on close comparisons and not of the inference
procedure; no query budget can remove it, since Section~\ref{sec:density} shows
that even exhaustive querying leaves a floor of the same character.

\paragraph{Concentration of error at the top of the order.}
At $n = 500$ we measure precision of $1/10$, $9/25$ and $24/50$ at
$k = 10, 25, 50$. Reference ties account for part of this --- $13$ items lie at
or above the tenth-highest label, so $X_{10}^{*}$ is not unique --- but the
dominant cause is visible in Table~\ref{tab:gap}: the top ${\sim}50$ items are
separated by gaps below $0.05$, where accuracy is $53.7\%$. The procedure
localizes the correct region, top-$50$ recall of $48\%$ exceeding chance by
$4.8\times$ and top-$25$ recall of $36\%$ by $7\times$, while being unable to
order within it. Consistent with Section~\ref{sec:topk}, precision is the
harshest available reading of this outcome: it charges full cost for exchanging
items of nearly identical label. Table~\ref{tab:sstregret} evaluates the same
output under both measures of Eq.~\eqref{eq:regret}.

\begin{table}[h]
\centering
\begin{tabular}{lrrrrr}
\toprule
$k$ & Precision & Tie-tolerant & $\bar q(S)$ & $\bar q(X_k^{*})$ & Regret \\
\midrule
$5$   & $1/5$    & $1/5$    & 0.8917 & 0.9861 & 0.2069 \\
$10$  & $2/10$   & $2/10$   & 0.8694 & 0.9736 & 0.2346 \\
$25$  & $11/25$  & $13/25$  & 0.8556 & 0.9433 & 0.2121 \\
$50$  & $26/50$  & $26/50$  & 0.8506 & 0.9067 & 0.1488 \\
$100$ & $62/100$ & $62/100$ & 0.8058 & 0.8612 & 0.1671 \\
\bottomrule
\end{tabular}
\caption{Selection quality on the $n=500$ SST instance under both measures.
The reference list mean is $\bar q(X) = 0.5295$.}
\label{tab:sstregret}
\end{table}

The divergence anticipated in Section~\ref{sec:topk} is confirmed on real data
and is large. At $k = 10$, precision of $0.2$ coincides with regret of $0.235$:
the returned subset attains mean label $0.869$ against an optimal $0.974$ and a
list mean of $0.530$, thus capturing ${\sim}77\%$ of the available quality margin.
Admitting label ties changes precision by at most two items, confirming that
non-uniqueness of $X_k^{*}$ is not the explanation; the explanation is that
precision and regret weight near-tie exchanges incomparably. A procedure that
appears to fail by an order of magnitude under exact rank recovery is therefore
substantially effective under a subset-quality criterion, and reporting only the
former would materially misstate its behaviour.

\paragraph{Replication.}
The $n=500$ instance was evaluated twice under identical configuration, differing
only in the randomized presentation order of each query. The two runs used
$4227$ and $4241$ queries --- the schedule is data-dependent, since verdicts
determine subsequent merge paths --- and yielded $\tau_b$ of $0.6127$ and
$0.6096$, with the accuracy profile of Table~\ref{tab:gap} agreeing to within
$1.0$ percentage point in every stratum. Variability induced by presentation
symmetrization is thus small relative to the effects reported here.

\section{Discussion}

The results delineate which parts of the problem admit improvement by
computation and which do not. Query density improves rank recovery with strongly
diminishing returns and cannot pass an oracle-determined floor
(Section~\ref{sec:density}); adaptive targeting recovers only a small fraction
of the remaining margin. Conversely, the choice of objective has first-order
consequences: the same oracle and estimator yield either a poor or a good result
on the identical output depending on whether exact rank recovery or subset
quality is demanded, and the schedule that is efficient for one is wasteful for
the other. Where near-ties are intrinsic to the data, reporting a strict total
order asserts distinctions the evidence does not support; a partition into
equivalence classes is the more faithful summary, and identifying such classes
requires uncertainty quantification on $s$ rather than merely its point estimate.

\section{Limitations}
\label{sec:limitations}

\begin{itemize}
  \item \textbf{Criterion--label mismatch.} The natural-language criterion is a
    proxy for the SST annotation protocol; discordance attributable to this
    mismatch is not separated from oracle error, so the accuracy ceiling of
    Table~\ref{tab:gap} is an upper bound on the true error rate of the oracle
    with respect to its own criterion.
  \item \textbf{Single oracle instance.} All empirical results use one model on
    one sampled instance per size. The $n=500$ instance was replicated once,
    bounding the variability due to presentation symmetrization, but variability
    across item samples, criteria and oracles is not characterized, and no
    confidence intervals are reported.
  \item \textbf{The uncertainty study is exploratory.} Table~\ref{tab:conf}
    rests on three pairs and six queries per estimator. It supports a claim of
    difficulty discrimination, not of calibration; in particular the reading of
    score differences as log-odds is model-implied and not validated against
    observed frequencies.
  \item \textbf{Regret depends on an assumed utility.} Eq.~\eqref{eq:regret}
    equates utility with the mean ground-truth label of the selected subset,
    which is linear and indifferent to the dispersion within $S$. An application
    whose utility is concave, or which requires a guaranteed minimum quality per
    element, would rank the schedules of Table~\ref{tab:topk} differently.
  \item \textbf{Simulation is well-specified.} The synthetic oracle obeys the
    same parametric family used for inference, so the simulation does not probe
    misspecification. Furthermore the exhaustive row of Table~\ref{tab:density}
    duplicates merge-sort edges, over-weighting them.
  \item \textbf{Interpretation of magnitudes.} The fitted $s_i$ derive from
    graph position and the prior $\alpha$; they are not estimates of a semantic
    distance and are comparable only within a single instance.
\end{itemize}

\section{Directions}

\begin{enumerate}
  \item \textbf{Objective-directed scheduling.} Section~\ref{sec:topk} indicates
    that the schedule should follow from the declared objective --- full order,
    top-$k$ subset, or partition into equivalence classes --- rather than
    defaulting to a total order.
  \item \textbf{Span-efficient schedules.} Proposition~\ref{prop:span} shows
    merge sort is asymptotically suboptimal in span. Round-based and
    sorting-network schedules attain $O(r)$ and $O(\log^2 n)$ span respectively;
    whether a round-based bootstrap matches merge sort in rank recovery at equal
    budget is open, and is suggested by the small margin of targeting over random
    selection in Table~\ref{tab:density}.
  \item \textbf{Uncertainty-directed querying.} Prioritizing pairs by the
    log-probability estimator of Section~\ref{sec:conf}, and admitting soft edge
    weights into Eq.~\eqref{eq:mm} only once calibration is established, since
    uncalibrated weights would allow confidently erroneous edges to dominate the
    likelihood.
  \item \textbf{Equivalence-class recovery.} Estimating uncertainty on $s$, by
    resampling the edge set or from the observed information, so that
    statistically indistinguishable items can be reported as tied.
  \item \textbf{Budget allocation near a boundary.} For a top-$k$ objective,
    concentrating queries at the rank-$k$ boundary rather than uniformly.
  \item \textbf{Calibration measurement.} Establishing, per oracle, the map from
    fitted score differences to realized win frequencies, which is prerequisite
    to the preceding items.
\end{enumerate}

\begin{thebibliography}{99}
\bibitem{bradley1952} R.~A. Bradley and M.~E. Terry.
  Rank analysis of incomplete block designs: I. The method of paired
  comparisons. \emph{Biometrika}, 39(3/4):324--345, 1952.
\bibitem{zermelo1929} E.~Zermelo.
  Die Berechnung der Turnier-Ergebnisse als ein Maximumproblem der
  Wahrscheinlichkeitsrechnung. \emph{Mathematische Zeitschrift},
  29:436--460, 1929.
\bibitem{hunter2004} D.~R. Hunter.
  MM algorithms for generalized Bradley--Terry models.
  \emph{Annals of Statistics}, 32(1):384--406, 2004.
\bibitem{negahban2017} S.~Negahban, S.~Oh, and D.~Shah.
  Rank Centrality: ranking from pairwise comparisons.
  \emph{Operations Research}, 65(1):266--287, 2017.
\bibitem{kemeny1959} J.~G. Kemeny.
  Mathematics without numbers. \emph{Daedalus}, 88(4):577--591, 1959.
\bibitem{dwork2001} C.~Dwork, R.~Kumar, M.~Naor, and D.~Sivakumar.
  Rank aggregation methods for the web. In \emph{WWW}, 2001.
\bibitem{kemenyElicit} Eliciting Kemeny rankings.
  arXiv:2312.11663. \url{https://arxiv.org/html/2312.11663v1}
\bibitem{activeRanking} Robust active ranking from sparse noisy comparisons.
  arXiv:1502.05556. \url{https://arxiv.org/abs/1502.05556v1}
\bibitem{zheng2023} L.~Zheng et al.
  Judging LLM-as-a-judge with MT-Bench and Chatbot Arena.
  In \emph{NeurIPS}, 2023.
\bibitem{tian2023} K.~Tian et al.
  Just ask for calibration: strategies for eliciting calibrated confidence
  scores from language models. In \emph{EMNLP}, 2023.
\bibitem{socher2013} R.~Socher et al.
  Recursive deep models for semantic compositionality over a sentiment
  treebank. In \emph{EMNLP}, 2013.
  \url{https://nlp.stanford.edu/sentiment/}
\bibitem{kendall1945} M.~G. Kendall.
  The treatment of ties in ranking problems.
  \emph{Biometrika}, 33(3):239--251, 1945.
\end{thebibliography}

\end{document}
